% =====================================================================
% Plantilla de examen con la clase exam.
% - Versión del alumno: compilar tal cual.
% - Clave del profesor: quitar el % de la línea \printanswers.
% - Versión B: copiar el archivo y cambiar \newcommand{\version}{A}
%   y los datos numéricos.
% =====================================================================
\documentclass[11pt,addpoints]{exam}
\usepackage[T1]{fontenc}
\usepackage[utf8]{inputenc}
\usepackage[spanish,es-noshorthands]{babel}
\usepackage{amsmath,amssymb}
\usepackage[margin=2.2cm]{geometry}
\usepackage{xcolor}

% \printanswers
\newcommand{\version}{A}
\renewcommand{\solutiontitle}{\noindent\textbf{Solución:}\enspace}
\shadedsolutions
\definecolor{SolutionColor}{rgb}{0.93,0.96,1}

\pagestyle{headandfoot}
\firstpageheader{Universidad Panamericana}{}{Versión \version}
\runningheader{Finanzas I}{Examen parcial 1}{Versión \version}
\footer{}{Página \thepage\ de \numpages}{}
\pointpoints{punto}{puntos}
\hqword{Pregunta}\hpword{Puntos}\hsword{Obtenidos}\htword{Total}

\begin{document}

\begin{center}
{\Large\bfseries Finanzas I --- Examen parcial 1}\\[4pt]
Duración: 90 minutos \quad\textbullet\quad Se permite calculadora financiera
\end{center}

\vspace{6pt}
\noindent Nombre: \makebox[0.55\textwidth]{\hrulefill}\quad Grupo: \makebox[0.15\textwidth]{\hrulefill}

\begin{center}
\gradetable[h][questions]
\end{center}

\begin{questions}

\question[10] Calcula el valor presente de \$10{,}000 que se recibirán dentro de 3 años si la tasa de descuento anual es $r = 10\%$.
\begin{solution}
$VP = \dfrac{10\,000}{(1.10)^3} = \dfrac{10\,000}{1.331} = 7\,513.15$.
\end{solution}

\question Un proyecto requiere una inversión inicial de \$1{,}000 y genera flujos de \$600 al final de los años 1 y 2. Con una tasa de descuento de $8\%$:
\begin{parts}
\part[5] Calcula el VPN.
\begin{solution}
$VPN = -1000 + \dfrac{600}{1.08} + \dfrac{600}{1.08^2} = -1000 + 555.56 + 514.40 = 69.96$.
\end{solution}
\part[5] ¿Conviene aceptar el proyecto? Justifica.
\begin{solution}
Sí: el VPN es positivo, de modo que el proyecto crea valor a la tasa de descuento dada.
\end{solution}
\end{parts}

\question[5] La tasa interna de retorno (TIR) de un proyecto es la tasa que:
\begin{choices}
  \choice maximiza el VPN.
  \CorrectChoice hace que el VPN sea igual a cero.
  \choice iguala la tasa libre de riesgo.
  \choice minimiza el periodo de recuperación.
\end{choices}

\question[5] Explica en tres renglones por qué un peso hoy vale más que un peso dentro de un año.
\begin{solution}
Porque el peso de hoy puede invertirse y generar rendimiento; además, el futuro conlleva incertidumbre e inflación. La tasa de descuento compensa ambos efectos.
\end{solution}
\makeemptybox{3cm}

\end{questions}

\end{document}
